\documentclass[11pt]{article}
\usepackage[letterpaper,margin=0.5in]{geometry}
\usepackage[T1]{fontenc}
\usepackage{lmodern}
\usepackage{microtype}
\usepackage{enumitem}
\usepackage{titlesec}
\usepackage[hidelinks]{hyperref}
\input{glyphtounicode}
\pdfgentounicode=1

\pagestyle{empty}
\setlength{\parindent}{0pt}
\setlength{\parskip}{0pt}
\setlist[itemize]{leftmargin=15pt,label=\textbullet,itemsep=1pt,topsep=1pt,parsep=0pt,partopsep=0pt}
\titleformat{\section}{\large\scshape}{}{0pt}{}[\titlerule]
\titlespacing*{\section}{0pt}{4pt}{2pt}
\urlstyle{same}
\hypersetup{pdftitle={Daniel David -- Resume},pdfauthor={Daniel David}}

\newcommand{\entry}[2]{\textbf{#1}\hfill #2\par}
\newcommand{\project}[2]{\textbf{#1} {\small\textit{\textbar\ #2}}\par}

\begin{document}
\begin{center}
    {\Huge\scshape Daniel David}\\[3pt]
    \small +7 (923) 408 72 93 \enspace\textbar\enspace
    \href{mailto:dayveed.d@gmail.com}{\underline{dayveed.d@gmail.com}} \enspace\textbar\enspace
    \href{https://www.linkedin.com/in/dayveedaniel}{\underline{LinkedIn}} \enspace\textbar\enspace
    \href{https://github.com/dayveedaniel}{\underline{github.com/dayveedaniel}}\\[3pt]
    \textbf{Spoken Languages:} English, Russian
\end{center}
\vspace{-9pt}
\fontsize{10}{11.5}\selectfont

\section{Education}
\entry{Innopolis University}{Innopolis, Russia}
\textit{MSc, Artificial Intelligence \& Data Engineering}\hfill Aug. 2024 -- Jul. 2026 \textbar\ CGPA: 4.24\par
\vspace{3pt}
\entry{Tomsk Polytechnic University}{Tomsk, Russia}
\textit{BEng, Software Engineering}\hfill Sep. 2019 -- Jul. 2024 \textbar\ CGPA: 3.85\par

\section{Experience}
\entry{PRAAMS --- Senior Software Developer}{Sep. 2022 -- Present}
\begin{itemize}
    \item Led architecture and development of PRAAMS, an AI-powered co-investor app combining conversational AI, stock analysis, and interactive financial visualizations.
    \item Conducted technical interviews and onboarded new team members.
\end{itemize}
\vspace{3pt}
\entry{Bellissimo Pizza --- ML / Data Engineering Consultant}{Jun. 2026 -- Nov. 2026}
\begin{itemize}
    \item Built an end-to-end data pipeline for order demand and price prediction using Airflow, PostgreSQL, and MinIO; automated ingestion, transformation, and training-data preparation.
    \item Contributed to hourly demand forecasting for 50+ restaurants over 24/48/72-hour horizons, incorporating POS and weather features into model training and evaluation.
    \item Automated model retraining on new data and maintained data pipelines to support recurring forecasting workflows.
\end{itemize}
\vspace{3pt}
\entry{Innopolis University --- Student Researcher}{Aug. 2025 -- Sep. 2026}
\begin{itemize}
    \item Implemented and trained ML/RL models in PyTorch; built web-scraping and big-data analysis pipelines using MongoDB, PySpark, and Apache Hadoop.
\end{itemize}

\section{Selected Projects}
\project{nanoTitan --- Open-Source Mixture-of-Experts Training Stack}{Python, PyTorch, Distributed Training, SDPA}
\begin{itemize}
    \item Added configurable SDPA / Flash Attention support to the existing multi-head attention module, retaining an eager attention path as a fallback.
    \item Integrated fused QKV projections and RMSNorm into the model implementation.
\end{itemize}
\vspace{3pt}
\project{Multilingual Math Reasoning for Nigerian Languages}{Docker, PostgreSQL, Hugging Face, vLLM, MLX}
\begin{itemize}
    \item Built a pipeline translating GSM8K math problems and solutions into Igbo, Hausa, and Yoruba for Tiny Aya supervised fine-tuning; produced the \href{https://huggingface.co/datasets/dayveedd/gsm8k-lowres-math}{\underline{gsm8k-lowres-math dataset}}.
    \item Benchmarked Qwen2.5-1.5B-Instruct and Tiny Aya variants on math tasks in all three languages to select a fine-tuning base model.
\end{itemize}
\vspace{3pt}
\project{ECG Benchmarking for Atrial Fibrillation Detection}{PySpark, PostgreSQL, MinIO}
\begin{itemize}
    \item Contributed to a reproducible benchmarking framework spanning 42+ ECG databases, with deterministic data splits, cached feature generation, and traceable experiment artifacts.
    \item Helped benchmark 14+ detectors across datasets using window- and recording-level metrics, with quality and runtime comparisons through configurable experiments.
\end{itemize}

\vspace{3pt}
\project{RAG for Document Collections}{Python, SentenceTransformers, LLMs}
\begin{itemize}
    \item Contributed to a Wikipedia-based RAG pipeline with article cleaning, chunking, 384-dimensional MiniLM embeddings, and incremental embedding updates.
    \item Combined 3-gram candidate retrieval with cosine-similarity reranking to supply relevant document snippets for LLM-generated answers.
\end{itemize}

\section{Research}
\textbf{Meta-learners and LLM-Guided Adaptive Weighting for Physics-Informed Neural Networks}\par
\begin{itemize}
    \item Lead-authored a study introducing an LLM meta-controller that dynamically balances PDE-residual and boundary-condition losses using training history, validation trends, and stability diagnostics.
    \item Benchmarked LLM-guided weighting against fixed weights and LBPINN on Poiseuille flow across 100 pressure-gradient values and five random seeds, assessing convergence, solution accuracy, and adaptation trade-offs.
\end{itemize}

\vspace{3pt}
\textbf{Software Engineering Metrics Selection Using Equilibrium Optimizer and Differential Evolution}\par
\begin{itemize}
    \item Implemented the Equilibrium Optimizer for software-metric selection, using heuristic search to identify compact feature subsets while preserving geometric relationships between data points.
    \item Used Sammon error and Kruskal stress to evaluate structural preservation; the Equilibrium Optimizer reduced 30 metrics to five in a co-authored study also evaluating Differential Evolution.
\end{itemize}

\end{document}
